\chapter{Plotting with JarvisPLOT}
\label{ch:jarvisplot}

\newthought{Figures are a plugin concern too.} \JPlot\ is the \textsc{yaml}-driven
plotting engine of the \Jarvis\ family: you describe a figure --- data sources, layers,
axes, styling --- in a \emph{scene} file, and the engine renders it. \Jtwo\ bridges it
as \cli{Jarvis2 plot}, so upgrading \JPlot\ brings new figure types without a runtime
release. This chapter covers the bridge and the scan-output hooks; the scene language
itself is \JPlot's own manual.

\section{Installing and invoking}
\label{sec:plot-install}

\begin{terminal}
\begin{jcmd}
python3 -m pip install -e '.[distributed,plot]'   # the plot extra

Jarvis2 plot scenes/money_plot.yaml               # render a scene
Jarvis2 plot scenes/money_plot.yaml               # re-render after edits
\end{jcmd}
\end{terminal}

The bridge is deliberately thin: it imports the installed engine, hands over the scene
path, and maps the result to exit codes --- \code{0} rendered, \code{1} plot failure,
\code{2} usage or missing package (with the \code{pip install} hint). The legacy
spelling \cli{Jarvis2 <scene>.yaml --plot} still routes here with a deprecation
warning.

\begin{jwarn}
The argument is a \JPlot\ \emph{scene}, not a scan task card. Handing \cli{plot} your
scan \textsc{yaml} fails cleanly --- scan-driven scene generation (``plot this run'')
is a planned convenience, not a current one; today you write the scene and point it at
the outputs below.
\end{jwarn}

\section{Feeding scenes from scan outputs}
\label{sec:plot-data}

Everything a scene needs is in the stable output tree (Chapter~\ref{ch:outputs}):

\begin{itemize}
  \item \file{DATABASE/samples.csv} --- the universal data source: full-column,
        flat, immediately consumable by \JPlot's tabular readers (and every scatter,
        histogram, and profile layer). Column selections and derived quantities are
        scene-side expressions over these columns.
  \item \file{DATABASE/samples.hdf5} --- the same records for readers that prefer
        \textsc{hdf5}.
  \item \file{levelset.json} --- \sampler{AdaptiveBridson} runs: crossing points and
        polylines, with stock overlay scenes shipped by \JPlot\ for drawing the traced
        contour over a sample scatter (Chapter~\ref{ch:adaptive-levelset}).
  \item \file{images/<scan>/flowchart.json} --- the semantic workflow flowchart
        exported by every run (Section~\ref{sec:flowchart}); \JPlot\ renders it to a
        publishable diagram.
\end{itemize}

\section{A working pattern}
\label{sec:plot-pattern}

Keep scenes in the project (\file{scenes/}), reference outputs with relative paths,
and treat figures as build artifacts:

\begin{terminal}
\begin{jcmd}
Jarvis2 run bin/scan.yaml           # produce outputs/scan-01/...
Jarvis2 plot scenes/overview.yaml   # scatter + LogL colormap
Jarvis2 plot scenes/contour.yaml    # levelset overlay
\end{jcmd}
\end{terminal}

Because \file{samples.csv} is contract-stable, scenes survive re-runs and framework
upgrades; because scenes are \textsc{yaml} in your project, \cli{project pack --share}
ships the figures' full provenance alongside the data
(Chapter~\ref{ch:project-cli}).

\begin{jtip}
For exploratory work, nothing beats \code{pandas} on \file{samples.csv} in a notebook.
Reach for \JPlot\ when a figure graduates to \emph{repeatable} --- the paper plot you
will regenerate twelve times as the scan grows.
\end{jtip}
